\section{Objetivos del Proyecto}

Los objetivos del presente estudio se han estructurado en dos categor\'ias. Los objetivos t\'ecnicos definen las metas cuantificables del an\'alisis de ingenier\'ia, mientras que los objetivos estrat\'egicos enmarcan el aporte del estudio a la gesti\'on de activos de Ingredion~S.A. La Tabla~\ref{tab:objetivos} presenta ambas categor\'ias de forma detallada.

\begin{table}[H]
	\centering
	\caption{Objetivos t\'ecnicos y estrat\'egicos del estudio.}
	\label{tab:objetivos}
	\setcounter{itemcount}{0}
	\begin{tabularx}{\textwidth}{@{}NlX@{}}
		\toprule
		\multicolumn{1}{c}{\textbf{Item}} & \textbf{Categor\'ia} & \textbf{Descripci\'on} \\
		\midrule
		\multicolumn{3}{l}{\textbf{Objetivos T\'ecnicos}} \\
		\cmidrule(lr){1-3}
		& Coeficiente $U$ & Determinar de manera independiente el coeficiente global de transferencia de calor $U$ [W/m\textsuperscript{2}$\cdot$°C] del sistema de calentamiento por media ca\~na rectangular en espiral, tanto en condici\'on est\'atica como durante la descarga de glucosa al carrotanque. \\
		& Perfiles de temperatura & Obtener los perfiles de temperatura de la glucosa Globe~1130 en funci\'on del tiempo y del nivel del tanque para tres escenarios de operaci\'on, evaluando el gradiente t\'ermico vertical en ausencia de agitaci\'on mec\'anica. \\
		& Din\'amica de descarga & Analizar la operaci\'on c\'iclica de carga a carrotanques de 24~toneladas con una frecuencia de 8~cargas en 24~horas, determinando los tiempos de calentamiento y el comportamiento t\'ermico entre descargas. \\
		& An\'alisis estructural & Analizar la integridad estructural del fondo toriesf\'erico (construido), el cuerpo cil\'indrico (existente) y la tapa c\'onica (existente) de acuerdo con las normas API~650 y ASME~Section~VIII, considerando presi\'on hidrost\'atica, viento, sismo y corrosi\'on. \\
		& Aislamiento t\'ermico & Determinar el espesor recomendado de aislamiento de lana mineral para la chaqueta de media ca\~na, con base en criterios de eficiencia energ\'etica y seguridad al contacto. \\
		& Metodolog\'ia de simulaci\'on & Establecer la metodolog\'ia detallada para la simulaci\'on CFD (COMSOL Multiphysics) del modelo t\'ermico y la simulaci\'on FEA (ANSYS) del modelo estructural. \\
		\midrule
		\multicolumn{3}{l}{\textbf{Objetivos Estrat\'egicos}} \\
		\cmidrule(lr){1-3}
		& An\'alisis independiente & Proporcionar a Ingredion~S.A. un an\'alisis t\'ecnico independiente de alta calidad sobre el comportamiento t\'ermico y la integridad mec\'anica del tanque de almacenamiento, como insumo para la toma de decisiones operativas. \\
		& Optimizaci\'on de proceso & Identificar las condiciones de operaci\'on \'optimas (temperatura del agua, nivel de llenado, espesor de aislamiento) para maximizar la eficiencia del sistema de calentamiento con la infraestructura existente. \\
		& Documentaci\'on t\'ecnica & Entregar a Ingredion~S.A. la documentaci\'on de soporte para la operaci\'on, inspecci\'on y mantenimiento del tanque, incluyendo recomendaciones fundamentadas en c\'alculos de ingenier\'ia y normativas internacionales. \\
		\bottomrule
	\end{tabularx}
\end{table}
